\subsection{Removing an odd relative defect for quadratic questions}

The distinction between ordinary relative characters and the cubic
crossed defect allows the following exact reduction.

\begin{lemma}\label{lem:odd-relative-kernel}
Let \(G=DH\) be a profinite group with \(D\) closed and normal,
and put \(K=D\cap H\). Suppose \(E\) is a closed normal subgroup
of \(G\), with \(K\subset E\subset D\), such that \(D/E\) is
pro-\(3\). For every finite abelian \(2\)-group \(A\), restriction
is a canonical bijection
\[
 \begin{aligned}
 \{\chi\in\operatorname{Hom}_{\rm cts}(G,A):\chi(H)=0\}
 &\longrightarrow\\
 &\{\psi\in\operatorname{Hom}_{\rm cts}(E,A)^G:\psi(K)=0\}.
 \end{aligned}
\]
Consequently the maximal pro-\(2\) quotient of the relative
abelian group can be computed as the maximal pro-\(2\) quotient of
\[
                  E/\overline{[G,E]K}.
\]
This does not assert \(G=EH\) or equality of the full relative
abelian quotients before taking their pro-\(2\) parts.
\end{lemma}
\begin{proof}
Let \(\psi\) be a character on the right and put \(L=\ker\psi\).
It is closed and normal in \(G\). In \(D/L\), the subgroup
\(E/L\) is a finite central \(2\)-group, and its quotient is
pro-\(3\). This group has a unique pro-\(3\) complement.
To see this directly, pass to finite quotients. A finite
central extension of a \(3\)-group by a \(2\)-group is the
direct product of its central \(2\)-subgroup with its Sylow
\(3\)-subgroup: a Sylow \(3\)-subgroup has the full required
order, and all its conjugates are itself because the
central subgroup and that Sylow subgroup generate the group.
The unique Sylow \(3\)-subgroups are compatible under quotient
maps, so their inverse limit is the asserted unique complement.
Hence
\[
                    D/L=(E/L)\times P
\]
with \(P\) pro-\(3\). Projection to \(E/L\), followed by
\(\psi\), extends \(\psi\) uniquely to a continuous
homomorphism \(\widetilde\psi:D\to A\).
Uniqueness follows also because a pro-\(3\) group has no
nontrivial continuous homomorphism to \(A\).

For every \(g\in G\), the conjugate of
\(\widetilde\psi\) is another extension of \(\psi\);
thus it is the same extension. The character on \(D\)
is therefore \(G\)-invariant and zero on \(K\).
It extends to \(G=DH\) by
\(\chi(dh)=\widetilde\psi(d)\).
The intersection and invariance conditions give
well-definedness and multiplicativity, and the compact
surjection \(D\times H\to G\) gives continuity.
Conversely every character on the left restricts as
asserted, and uniqueness on \(D\), followed by \(G=DH\),
proves injectivity. Taking all finite abelian \(2\)-groups
gives the pro-\(2\) assertion.
\end{proof}

Normality of $D$ is necessary. Let $G=C_2\times\mathfrak S_4$,
$D=C_2\times\langle(123)\rangle$, and $E=C_2\times\{1\}$.
In the dihedral subgroup $J=\langle r=(1234),s=(13)\rangle$,
let $\theta:J\to C_2$ have $\theta(r)=0$, $\theta(s)=1$,
and let $H$ be its graph in $G$. Then $G=DH$, $D\cap H=1$,
and $D/E=C_3$, while $E$ is central. The nonzero character of
$E$ cannot extend to a character of $G$ killing $H$:
every character of $G$ is a combination of the first projection
and the sign character, and $\theta$ is independent of the
restriction of the sign to $J$. Thus dropping normality would
invalidate the displayed bijection.

Apply this lemma to the radical--CM defect of
Section~\ref{sec:lr-relative-gate}
\[
 d:G\longrightarrow C,\qquad d|_H=0,\qquad
 d(gh)=d(g)+\lambda_gd(h).
\]
Let \(d_{(3)}\) be its order-three component and put
\[
                  D^{\rm cub}=\ker(d_{(3)}|_D).
\]
On \(D\), the crossed law is an ordinary homomorphism.
Its \(G\)-equivariance makes \(D^{\rm cub}\) normal in \(G\);
it contains \(K=D\cap H\), and \(D/D^{\rm cub}\) has
exponent three. Thus the ordinary relative quadratic
character problem, and more generally the maximal
pro-\(2\) relative abelian quotient, may be computed on
\(D^{\rm cub}\) without proving \(d_{(3)}=0\) on \(G\).
On this smaller domain the cubic radical--CM compatibility
is imposed by definition. The tame harmonic square tests
remain arithmetic-trivial quadratic characters there.
If the cubic defect is nonzero, this is a proper reduction;
if it vanishes, \(D^{\rm cub}=D\).

More precisely, use the established closed normal subgroup
\(R_D\subset D\) with \(G=R_DH\), and put
\[
 S=R_D\cap\ker d_{(3)},\qquad K_R=R_D\cap H.
\]
Then \(K_R\subset S\), the subgroup \(S\) is closed and normal
in \(G\), and \(R_D/S\) has exponent three. The preceding
lemma, applied with \(R_D\) in place of \(D\), gives the
canonical identification
\[
 B_{\mathrm{max}\,2}\ \simeq\
 \left(S/\overline{[G,S]K_R}\right)_{\mathrm{max}\,2},
\]
where the subscript denotes the maximal pro-\(2\) quotient.
For the original invariant-character group
\(X_2(N)=\operatorname{Hom}_{\rm cts}(N,\mathbf F_2)^G\),
restriction to \(S\) has exactly the same kernel as
restriction to \(R_D\). Indeed a character of \(N\) that is
zero on \(S\), when restricted to \(R_D\), factors through
the pro-\(3\) quotient \(R_D/S\) and is therefore zero.
The converse is immediate. Thus this extra cubic-kernel
restriction loses none of the original quadratic characters
that survived the \(R_D\) normalization.
It is a proper restriction precisely when
\(d_{(3)}|_{R_D}\ne0\); no such nonvanishing is asserted.
